\documentclass{ximera}

\input{../../../preamble.tex}


\definecolor{fvdnavy}{HTML}{071D43}
\definecolor{fvdcyan}{HTML}{20BFC6}
\definecolor{fvdcyanlight}{HTML}{EFFCFB}
\definecolor{fvdmagenta}{HTML}{F10D69}
\definecolor{fvdmagentalight}{HTML}{FFF1F7}
\definecolor{fvdtext}{HTML}{163044}





%=================================================
% Pedagogische omgevingen
% PDF: tcolorbox
% HTML: learning-box
%=================================================

\newenvironment{voorkennis}
{%
    \ifdefined\HCode
        \HCode{%
<section class="learning-box learning-box--prior">
<div class="learning-box__header">
<span class="learning-box__icon" aria-hidden="true"></span>
<span class="learning-box__title">Voorkennis</span>
</div>
<div class="learning-box__content">
        }%
        \begin{itemize}
    \else
       \begin{tcolorbox}[
    enhanced,
    breakable,
    colback=fvdcyanlight,
    colframe=fvdcyan,
    coltext=fvdtext,
    boxrule=0.5pt,
    leftrule=4pt,
    arc=4mm,
    outer arc=4mm,
    left=7mm,
    right=6mm,
    top=5mm,
    bottom=5mm,
    before skip=6mm,
    after skip=6mm,
    title=Voorkennis,
    coltitle=fvdcyan,
    fonttitle=\bfseries\Large,
    attach title to upper={\par\smallskip},
    boxed title style={
        empty,
        left=0mm,
        right=0mm,
        top=0mm,
        bottom=0mm
    },
    overlay={
        \draw[
            fvdcyan,
            line width=0.5pt
        ]
        ([xshift=42mm,yshift=-13mm]frame.north west)
        --
        ([xshift=-7mm,yshift=-13mm]frame.north east);
    }
]
        \begin{itemize}
    \fi
}
{%
    \end{itemize}
    \ifdefined\HCode
        \HCode{%
</div>
</section>
        }%
    \else
        \end{tcolorbox}
    \fi
}


\newenvironment{leerdoelen}
{%
    \ifdefined\HCode
        \HCode{%
<section class="learning-box learning-box--goals">
<div class="learning-box__header">
<span class="learning-box__icon" aria-hidden="true"></span>
<span class="learning-box__title">Leerdoelen</span>
</div>
<div class="learning-box__content">
        }%
        \begin{itemize}
    \else
\begin{tcolorbox}[
    enhanced,
    breakable,
    colback=fvdmagentalight,
    colframe=fvdmagenta,
    coltext=fvdtext,
    boxrule=0.5pt,
    leftrule=4pt,
    arc=4mm,
    outer arc=4mm,
    left=7mm,
    right=6mm,
    top=5mm,
    bottom=5mm,
    before skip=6mm,
    after skip=6mm,
    title=Leerdoelen,
    coltitle=fvdmagenta,
    fonttitle=\bfseries\Large,
    attach title to upper={\par\smallskip},
    boxed title style={
        empty,
        left=0mm,
        right=0mm,
        top=0mm,
        bottom=0mm
    },
    overlay={
        \draw[
            fvdmagenta,
            line width=0.5pt
        ]
        ([xshift=42mm,yshift=-13mm]frame.north west)
        --
        ([xshift=-7mm,yshift=-13mm]frame.north east);
    }
]
        \begin{itemize}
    \fi
}
{%
    \end{itemize}
    \ifdefined\HCode
        \HCode{%
</div>
</section>
        }%
    \else
        \end{tcolorbox}
    \fi
}


\addPrintStyle{../../..}

\title{Oefeningen --- Lineaire regressie}
\author{Ingmar Herreman}


\begin{document}

% FVD_CHAPTER_DATA_START
% Automatisch gegenereerd uit index.tex.
% Niet handmatig aanpassen.
\fvdchapterdata{3}{Verbanden onderzoeken met data}{3}
% FVD_CHAPTER_DATA_END












\maketitle




\begin{oefening}[Parameters interpreteren]{\oefB \ESS}

Voor een dataset werd het regressiemodel

\[
\hat y=-2{,}8x+74
\]

gevonden. $x$ wordt gemeten in ${}^{\circ}\mathrm C$ en $y$ in
$\mathrm{kWh}$.

\begin{question}
Wat is de eenheid van de richtingscoëfficiënt?

\begin{oplossing}
\[
\frac{\mathrm{kWh}}{{}^{\circ}\mathrm C}.
\]
\end{oplossing}
\end{question}

\begin{question}
Interpreteer $-2{,}8$.

\begin{oplossing}
Wanneer $x$ met $1\,{}^{\circ}\mathrm C$ toeneemt, daalt de door het
model voorspelde waarde van $y$ met $2{,}8\,\mathrm{kWh}$.
\end{oplossing}
\end{question}

\begin{question}
Wat betekent $74$ algebraïsch?

\begin{oplossing}
Voor $x=0\,{}^{\circ}\mathrm C$ voorspelt het model
$\hat y=74\,\mathrm{kWh}$.
\end{oplossing}
\end{question}
\end{oefening}


\begin{oefening}[Voorspellen met een model]{\oefB \ESS}

Voor een sensor geldt binnen het onderzochte bereik

\[
\hat U=0{,}42T+1{,}8,
\]

met $T$ in ${}^{\circ}\mathrm C$ en $U$ in volt.

\begin{question}
Voorspel $U$ bij $T=20\,{}^{\circ}\mathrm C$.

\begin{oplossing}
\[
\hat U=0{,}42\cdot20+1{,}8=10{,}2\,\mathrm V.
\]
\end{oplossing}
\end{question}

\begin{question}
Wat betekent het dakje in $\hat U$?

\begin{oplossing}
Het gaat om de door het regressiemodel voorspelde waarde, niet
noodzakelijk om de werkelijk gemeten waarde.
\end{oplossing}
\end{question}
\end{oefening}


\begin{oefening}[Interpolatie of extrapolatie?]{\oefB \ESS}

Een model is gebaseerd op metingen voor

\[
15\leq x\leq 80.
\]

Classificeer de volgende voorspellingen.

\[
x=20,\qquad x=57,\qquad x=80,\qquad x=95,\qquad x=4.
\]

\begin{oplossing}
$x=20$, $57$ en $80$ liggen binnen het gemeten bereik en zijn
interpolaties. $x=95$ en $x=4$ liggen buiten het bereik en zijn
extrapolaties.
\end{oplossing}
\end{oefening}


\begin{oefening}[Gemeten en voorspeld]{\oefB}

Voor een waarneming geldt

\[
y=53{,}7
\qquad\text{en}\qquad
\hat y=50{,}9.
\]

\begin{question}
Bereken het residu.

\begin{oplossing}
\[
e=y-\hat y=53{,}7-50{,}9=2{,}8.
\]
\end{oplossing}
\end{question}

\begin{question}
Ligt het meetpunt boven of onder de regressierechte?

\begin{oplossing}
Boven de regressierechte, want het residu is positief.
\end{oplossing}
\end{question}
\end{oefening}




\begin{oefening}[Groei van een jonge plant]{\oefU \ESS}

Van een jonge plant wordt gedurende de eerste weken de hoogte gemeten.

\[
\begin{array}{c|rrrrrrrr}
t\;(\mathrm{week})&1&2&3&4&5&6&7&8\\ \hline
h\;(\mathrm{cm})&8{,}2&11{,}1&14{,}5&17{,}2&20{,}6&23{,}1&26{,}5&29{,}4
\end{array}
\]

\begin{question}
Maak met ICT een spreidingsdiagram en beschrijf het patroon.
\end{question}

\begin{question}
Bepaal $r$ en voer een lineaire regressie uit.
\end{question}

\begin{question}
Noteer het model als

\[
\hat h=at+b.
\]
\end{question}

\begin{question}
Interpreteer $a$ met eenheid.

\begin{oplossing}
$a$ geeft de gemiddelde toename van de door het model voorspelde
planthoogte per week, in $\mathrm{cm/week}$.
\end{oplossing}
\end{question}

\begin{question}
Gebruik het model om de hoogte in week $6{,}5$ te schatten. Is dit
interpolatie of extrapolatie?

\begin{oplossing}
De voorspelling wordt verkregen door $t=6{,}5$ in het gevonden model
in te vullen. Het is interpolatie omdat $6{,}5$ binnen het gemeten
interval van $1$ tot $8$ weken ligt.
\end{oplossing}
\end{question}

\begin{question}
Is een voorspelling voor week $52$ verantwoord?

\begin{oplossing}
Niet zonder bijkomende informatie. Dat is een zeer verre extrapolatie.
Biologische groei blijft bovendien niet onbeperkt lineair.
\end{oplossing}
\end{question}
\end{oefening}


\begin{oefening}[Hoogte en temperatuur]{\oefU \ESS}

Op acht meetstations worden hoogte en gemiddelde temperatuur gemeten.

\[
\begin{array}{c|rrrrrrrr}
H\;(\mathrm m)&120&280&450&690&910&1180&1450&1710\\ \hline
T\;({}^{\circ}\mathrm C)&15{,}8&15{,}0&13{,}9&12{,}8&11{,}7&10{,}1&8{,}9&7{,}2
\end{array}
\]

\begin{question}
Maak een spreidingsdiagram en beoordeel of een lineair model zinvol
lijkt.
\end{question}

\begin{question}
Bepaal met ICT $r$ en de regressierechte.
\end{question}

\begin{question}
Interpreteer de richtingscoëfficiënt met eenheid.

\begin{oplossing}
De richtingscoëfficiënt geeft de verandering van de voorspelde
temperatuur per meter hoogte, in ${}^{\circ}\mathrm C/\mathrm m$.
Voor een praktischere interpretatie kan men de verandering per
$100\,\mathrm m$ berekenen.
\end{oplossing}
\end{question}

\begin{question}
Voorspel met het model de temperatuur op $1000\,\mathrm m$ hoogte en
classificeer de voorspelling.

\begin{oplossing}
Vul $H=1000$ in in het gevonden regressiemodel. Het is interpolatie,
want $1000\,\mathrm m$ ligt binnen het gemeten hoogtebereik.
\end{oplossing}
\end{question}

\begin{question}
Waarom moet je voorzichtig zijn met een voorspelling voor
$8000\,\mathrm m$?

\begin{oplossing}
Dat is een zeer verre extrapolatie. Het gevonden lokale lineaire
verband hoeft over zo'n groot hoogtebereik niet geldig te blijven.
\end{oplossing}
\end{question}
\end{oefening}


\begin{oefening}[Wanneer heeft het intercept betekenis?]{\oefU \ESS}

Een regressiemodel voor de massa $m$ van metalen staven in functie van
hun lengte $L$ is

\[
\hat m=2{,}31L+1{,}7.
\]

De gemeten lengtes liggen tussen $10\,\mathrm{cm}$ en
$50\,\mathrm{cm}$.

\begin{question}
Wat betekent de richtingscoëfficiënt?

\begin{oplossing}
Volgens het model neemt de voorspelde massa met ongeveer
$2{,}31\,\mathrm g$ toe per extra centimeter lengte.
\end{oplossing}
\end{question}

\begin{question}
Algebraïsch voorspelt het model bij $L=0$ een massa van
$1{,}7\,\mathrm g$. Waarom moet je die interpretatie kritisch bekijken?

\begin{oplossing}
$L=0$ ligt buiten het onderzochte bereik en een staaf met lengte nul
heeft in deze context geen gewone fysische betekenis. Het intercept
hoeft dus geen betekenisvolle fysische grootheid te zijn.
\end{oplossing}
\end{question}
\end{oefening}




\begin{oefening}[Kalibratie van een oplossing]{\oefU \ESS}

In een chemisch experiment wordt de absorbantie $A$ gemeten voor
oplossingen met bekende concentratie $c$.

\[
\begin{array}{c|rrrrrrrr}
c\;(\mathrm{mmol/L})&0{,}5&1{,}0&1{,}5&2{,}0&2{,}5&3{,}0&3{,}5&4{,}0\\ \hline
A&0{,}12&0{,}21&0{,}34&0{,}43&0{,}55&0{,}64&0{,}77&0{,}85
\end{array}
\]

Een onbekende oplossing heeft een gemeten absorbantie van $0{,}60$.

\begin{question}
Maak een spreidingsdiagram en beoordeel de lineariteit.
\end{question}

\begin{question}
Bepaal $r$ en een lineair regressiemodel voor $A$ in functie van $c$.
\end{question}

\begin{question}
Interpreteer de richtingscoëfficiënt.
\end{question}

\begin{question}
Gebruik het model om te onderzoeken welke concentratie ongeveer
overeenkomt met $A=0{,}60$.

\begin{oplossing}
Los in het gevonden model
\[
0{,}60=ac+b
\]
op naar $c$. De precieze waarde hangt af van de door ICT gevonden en
afgeronde regressiecoëfficiënten.
\end{oplossing}
\end{question}

\begin{question}
Leg uit waarom het meetbereik van de kalibratie belangrijk is voor de
betrouwbaarheid van deze schatting.

\begin{oplossing}
Binnen het gekalibreerde bereik hebben we gegevens die het lineaire
model ondersteunen. Buiten dat bereik zouden we extrapoleren en weten
we niet of dezelfde lineaire relatie geldig blijft.
\end{oplossing}
\end{question}
\end{oefening}


\begin{oefening}[Een model dat te ver wordt doorgetrokken]{\oefU \ESS}

Voor baby's tussen $1$ en $12$ maanden wordt een sterk lineair verband
gevonden tussen leeftijd $t$ in maanden en lichaamslengte $L$ in cm:

\[
\hat L=1{,}55t+52{,}4.
\]

\begin{question}
Interpreteer $1{,}55$.

\begin{oplossing}
Binnen het onderzochte leeftijdsbereik neemt de voorspelde
lichaamslengte volgens het model gemiddeld met ongeveer
$1{,}55\,\mathrm{cm}$ per extra maand toe.
\end{oplossing}
\end{question}

\begin{question}
Voorspel de lengte bij $t=8$ maanden. Is dit interpolatie of
extrapolatie?

\begin{oplossing}
\[
\hat L=1{,}55\cdot8+52{,}4=64{,}8\,\mathrm{cm}.
\]
Dit is interpolatie.
\end{oplossing}
\end{question}

\begin{question}
Bereken wat het model voor $t=240$ maanden voorspelt.

\begin{oplossing}
\[
\hat L=1{,}55\cdot240+52{,}4=424{,}4\,\mathrm{cm}.
\]
\end{oplossing}
\end{question}

\begin{question}
Wat leert dit resultaat over modelleren?

\begin{oplossing}
Een regressiemodel kan binnen het onderzochte bereik zeer bruikbaar
zijn en bij verre extrapolatie volkomen onrealistische voorspellingen
geven. Een model moet altijd in zijn context en toepassingsgebied
worden geïnterpreteerd.
\end{oplossing}
\end{question}
\end{oefening}




\begin{oefening}[Rechte of kromme?]{\oefG}

Een experiment levert de gegevens

\[
\begin{array}{c|rrrrrrrr}
x&1&2&3&4&5&6&7&8\\ \hline
y&1{,}1&3{,}9&9{,}2&16{,}1&24{,}8&36{,}3&49{,}1&64{,}2
\end{array}
\]

\begin{question}
Maak een spreidingsdiagram.
\end{question}

\begin{question}
Laat ICT toch een lineaire regressie uitvoeren.
\end{question}

\begin{question}
Leg uit waarom een regressierechte hier ondanks de software-uitvoer
geen goede modelkeuze is.

\begin{oplossing}
De punten volgen duidelijk een krom, ongeveer kwadratisch patroon.
De afwijkingen van een rechte zijn systematisch en niet willekeurig.
Een softwarepakket kan altijd een regressierechte berekenen, maar dat
betekent niet dat het model inhoudelijk geschikt is.
\end{oplossing}
\end{question}
\end{oefening}


\begin{oefening}[Onderzoek een eigen dataset]{\oefG}
  
Zoek of verzamel een kleine dataset met twee numerieke grootheden uit
fysica, chemie, biologie, technologie of een andere wetenschappelijke
context.

Voer een volledige analyse uit:

\begin{enumerate}
  \item beschrijf de grootheden, eenheden en bron van de gegevens;
  \item maak een spreidingsdiagram;
  \item analyseer richting, vorm, sterkte en bijzonderheden;
  \item bepaal $r$;
  \item beoordeel of een lineair model geschikt is;
  \item bepaal indien verantwoord de regressierechte;
  \item interpreteer de parameters;
  \item maak één zinvolle voorspelling;
  \item bepaal of dit interpolatie of extrapolatie is;
  \item formuleer minstens één beperking van je model.
\end{enumerate}

\begin{oplossing}
Er zijn veel correcte oplossingen mogelijk. De kwaliteit wordt bepaald
door de samenhang tussen dataset, grafische analyse, modelkeuze,
interpretatie en kritische terugkoppeling naar de context.
\end{oplossing}
\end{oefening}


\end{document}
